\begin{helper}[Terminal parallel-extension ordering case]
Let \(T\) be a positive-rung extended Tashkinov tree generated by
\[
 \mathcal T=\{(T_i,\varphi_{i-1},S_{i-1},F_{i-1},\Theta_{i-1}):
                   1\leq i\leq n+1\},
\]
and suppose that \(T\) has the maximum property. Keep the graph \(G\),
palette size \(k\), and distinguished uncolored edge \(e_0\) fixed, but
allow the certified input coloring to vary arbitrarily. Assume explicitly
the following six assertions for every positive-rung ETT \(T'\) on these
fixed data, with arbitrary certified input coloring \(\varphi'_0\), maximum
property, and rung number \(s<n\), writing \(\varphi'_s\) for its generating
coloring and \(\mathcal T'\) for its certified series:
\begin{enumerate}
\item \(V(T')\) is elementary with respect to \(\varphi'_s\).
\item Under every proper \((T'_{s+1},D'_s,\varphi'_s)\)-stable coloring
leaving precisely \(e_0\) uncolored, \(T'_{s+1}\) is interchangeable for
any pair of colors of which at least one is missing on \(T'_{s+1}\).
\item For each PE slot \(a\leq s\) with supporting vertex \(u_a\) and
\(m(u_a)=b\), every proper \((T'_a,D'_a,\varphi'_a)\)-stable coloring
leaving precisely \(e_0\) uncolored is
\((T'(u_a)-u_a,D'_{b-1},\varphi'_{b-1})\)-stable. When \(b>1\), it is
also \((T'_{b-1},D'_{b-1},\varphi'_{b-1})\)-stable. Moreover, the stored
color pairs at two distinct PE slots with distinct supporting vertices
and the same first-appearance index are disjoint.
\item If the terminal slot is PE with supporting vertex \(u_s\) and
stored colors \(\gamma'_s,\delta'_s\), then under every proper
\((T'_s,D'_s,\varphi'_s)\)-stable coloring leaving precisely \(e_0\)
uncolored, the \((\gamma'_s,\delta'_s)\)-component through \(u_s\)
contains exactly the one vertex \(u_s\) of \(T'_s\).
\item If \(\tau\) is any proper \((T'_s,D'_s,\varphi'_s)\)-stable
coloring leaving precisely \(e_0\) uncolored, \(d\) is defective on
\(T'_s\) under \(\tau\), and \(x\) is a nonminimum vertex of
\(I[\partial_{\tau,d}(T'_s)]\), then
\(x\preceq_{T'}u_a\) for every PE supporting vertex or SE extension
vertex \(u_a\) with \(m(x)\leq a\leq s\).
\item Every proper \((T'_s,D'_s,\varphi'_s)\)-stable coloring leaving
precisely \(e_0\) uncolored is a \(\varphi'_s\bmod T'_s\) coloring.
\end{enumerate}
Assume that the terminal positive slot of \(\mathcal T\) is PE.

Write \(T_n\) for the terminal pre-extension tree. Let \(\sigma_n\) be
any proper \((T_n,D_n,\varphi_n)\)-stable coloring leaving precisely
\(e_0\) uncolored, let \(d\) be defective on \(T_n\) under \(\sigma_n\),
and let \(v\) be a nonminimum vertex of
\(I[\partial_{\sigma_n,d}(T_n)]\). If \(v_i\) is the exact chosen
vertex at a PE or SE slot \(i\leq n\), and \(m(v)\leq i\), then
\[
                              v\preceq_T v_i .
\]
Here the weak order means equality or strict precedence in the exact
vertex order of the displayed final tree, and RE chosen vertices are not
included among the applicable vertices.
\end{helper}
\begin{proof}
Specializing the six lower-rung assumptions to the current certified
input gives the lower-rung premise of \noderef{ExtensionRules}, so current
assertion (iii) is available with all of its stability and disjointness
clauses. Applying \noderef{ExtensionBase} to the original generalized
lower-rung package gives current assertion (iv): under any terminal stable
coloring, the component of the terminal PE colors through its supporting
vertex meets \(T_n\) only at that vertex.

Fix \(\sigma_n,d,v,i,v_i\) as in the statement and put \(X=V(T_n)\).
By \noderef{BoundaryColorVertex}, choose an exact \(d\)-colored boundary
edge \(a\) incident with \(v\), and choose an earlier boundary vertex
\(w\prec_{T_n}v\) together with an exact \(d\)-colored boundary edge
\(b\) incident with \(w\). These are distinct exact edge identities:
\noderef{BoundaryEdges} gives one inside and one outside endpoint for a
boundary edge, while \(w\ne v\) follows from \noderef{TreeVertexOrder};
one boundary identity cannot consequently be incident inside \(X\) with
both vertices. Their membership and color equations are precisely those
of \noderef{BoundaryColorEdges}.

Unfold the terminal PE constructor of \noderef{TashkinovSeriesStep}.
Write \(v_n,f_n,\delta_n\) for its maximum-defective choice, write
\(\pi_{n-1}\) for the selected coloring, \(\pi'_{n-1}\) for the second
stable coloring, and \(\gamma_n\) for its PE color. The constructor
identifies the chosen vertex and stored colors with
\(v_n,\gamma_n,\delta_n\), and it identifies
\[
 \varphi_n=\pi'_{n-1}/P_{v_n}(\gamma_n,\delta_n,\pi'_{n-1}).
\]
The fields of \noderef{MaximumDefectiveChoice} say that
\(\pi_{n-1}\) is proper, leaves exactly \(e_0\) uncolored, is
\((T_n,D_{n-1},\varphi_{n-1})\)-stable, and gives the exact incident
edge \(f_n\) color \(\delta_n\). The PE constructor gives the
corresponding enlarged-control stability of \(\pi'_{n-1}\).

By \noderef{ParallelExtensionColor}, \(\gamma_n\) is missing at \(v_n\)
under \(\pi_{n-1}\), and the vertex clause of \noderef{StableColoring}
transfers that missingness to \(\pi'_{n-1}\). Its edge clause fixes the
\(\delta_n\)-colored identity \(f_n\). Properness and
\noderef{AdjacentEdges} make \(f_n\) the unique incident
\(\delta_n\)-edge. Therefore the displayed component switch changes
\(f_n\) from \(\delta_n\) to \(\gamma_n\), introduces no other incident
\(\delta_n\)-edge, and makes \(\delta_n\) missing at \(v_n\) under
\(\varphi_n\). Terminal stability transfers this missingness to
\(\sigma_n\).

The terminal part of \noderef{HkuParallelBoundaryStructure} gives the
identity equality
\[
                 \partial_{\varphi_n,\gamma_n}(X)=\{f_n\}.
\]
The terminal stored-set equation and
\noderef{TashkinovStoredColors} put \(\gamma_n\) in the terminal
missing-or-controlled set. Hence the forward edge clause of terminal
stability fixes every edge in that singleton. Conversely, the symmetric
criterion \noderef{sc1}, applied to the two proper exactly-one-uncolored
colorings, fixes every \(\sigma_n\)-edge of color \(\gamma_n\) back to
its \(\varphi_n\)-color. Since endpoints are unchanged,
\[
                 \partial_{\sigma_n,\gamma_n}(X)=\{f_n\}.       \tag{1}
\]
The boundary of the defective color \(d\) has at least two exact
identities by \noderef{DefectiveColor}; therefore \(d\ne\gamma_n\).

Current assertion (iv), obtained above from \noderef{ExtensionBase},
says that
\[
 K=P_{v_n}(\gamma_n,\delta_n,\sigma_n)
 \quad\hbox{satisfies}\quad K\cap X=\{v_n\}.                    \tag{2}
\]
Apply \noderef{ParallelExtensionUndoStable} and define
\(\tau=\sigma_n/K\). It yields properness, exact one-edge-uncoloredness,
and
\[
       \tau\text{ is }(T_n,D_{n-1},\varphi_{n-1})\text{-stable}. \tag{3}
\]

We check the required boundary transport edge by edge. Let \(g\) be an
exact identity in \(\partial_{\sigma_n,d}(X)\), and let \(x\in X\) be
its inside endpoint, supplied by \noderef{BoundaryEdges}. If the
edge-node \(g\) lay in \(K\), the incidence definition
\noderef{TwoColorComponent} would also place \(x\) in \(K\), so (2)
would give \(x=v_n\). Since \(g\) has color \(d\) in the two-color
component, its color is \(\gamma_n\) or \(\delta_n\). The first
alternative contradicts \(d\ne\gamma_n\), using (1); the second
contradicts the \(\delta_n\)-missingness at \(v_n\), expanded with
\noderef{MissingColors}. Thus no \(d\)-colored boundary identity is
switched. By \noderef{TwoColorComponentSwap}, its color and endpoints
are unchanged, so
\[
 \partial_{\sigma_n,d}(X)\subseteq\partial_{\tau,d}(X).         \tag{4}
\]
In particular the two distinct identities \(a,b\), their inside
endpoints \(v,w\), and the order \(w\prec_{T_n}v\) survive. Expanding
\noderef{DefectiveEdge} and \noderef{DefectiveVertex}, \(v\) is a
nonminimum defective vertex under \(\tau\).

If \(i=n\), the exact terminal kind and the applicability disjunction
force its PE alternative, and the chosen-vertex equation gives
\(v_i=v_n\). The maximality field of
\noderef{MaximumDefectiveChoice}, applied using (3), gives
\(v=v_n\) or \(v\prec_{T_n}v_n\). The ordered inclusion of \(T_n\) in
\(T\), from \noderef{ExtendedTashkinovTree} and
\noderef{TreeSequenceInitialSegment}, transports the strict alternative,
so \(v\preceq_Tv_i\).

Suppose finally that \(i<n\). The inequality \(m(v)\leq i\), the least
stage definition \noderef{ETTFirstAppearanceIndex}, and the nested
series stages put \(v\) in \(T_{n-1}\). Entrywise prefix equality and
vertex injectivity from \noderef{TreeSequence} put the earlier vertex
\(w\) there as well and retain \(w\prec_{T_{n-1}}v\). The outside
endpoints of \(a,b\) lie outside \(X\), hence outside the smaller tree.
Thus the same two exact identities show that \(d\) is defective on
\(T_{n-1}\) under \(\tau\), with \(v\) nonminimum.

Now \(1\leq i<n\), so \(n-1\) is positive.
\noderef{CertifiedETTRestriction} at rung \(n-1\) supplies an actual
positive-rung maximum ETT whose displayed tree is the old \(T_n\), whose
generating coloring is \(\varphi_{n-1}\), and whose retained trees,
stored sets, chosen vertices, kinds, and step certificates are exactly
those of the original series. Relation (3) restricts to its terminal
ladder tree \(T_{n-1}\): vertex missingness restricts directly, and
each controlled edge color has the same witnessing missing vertex in
\(T_n\), so the edge clause of \noderef{StableColoring} restricts as
well. The first-appearance minimum for \(v\) is unchanged because its
least stage is at most \(i<n\) and all candidate stages through \(i\)
are identified exactly. Apply lower-rung assertion (v) to this
restriction. It gives \(v=v_i\) or \(v\prec_{T_n}v_i\); the ordered
inclusion of \(T_n\) in \(T\) again gives \(v\preceq_Tv_i\).
\end{proof}

